\section{Componentes para escalabilidad y tolerancia a fallos}
\label{sec:escalabilidad}

La arquitectura base tiene una sola instancia de cada componente: es simple, pero
el gestor, el servicio de cancelación y persistencia son puntos únicos de fallo, y
el procesamiento de un único procesador RMC limita el número de transacciones por
minuto. Para la segunda entrega se incorporarán los componentes de la
Tabla~\ref{tab:componentes-ft}. Aquí solo se mencionan y justifican; los diagramas
completos se presentarán en esa entrega.

\begin{longtable}{p{3.6cm}p{4.4cm}p{7.3cm}}
\caption{Componentes de escalabilidad y tolerancia a fallos.}
\label{tab:componentes-ft}\\
\toprule
Componente & Qué resuelve & Justificación \\
\midrule
\endfirsthead
\toprule
Componente & Qué resuelve & Justificación \\
\midrule
\endhead
\textbf{Réplicas del procesador RMC} &
Escalabilidad horizontal de reservas, modificaciones y consultas. &
Los procesadores no guardan estado, así que cualquier réplica atiende cualquier
petición. El gestor ya reparte la carga entre los procesadores libres, de modo que
agregar réplicas no cambia el código. Es la diferencia principal entre la
arquitectura base y la escalable en las pruebas de la Tabla~1. \\
\textbf{Clientes \emph{Lazy Pirate}} &
Peticiones o respuestas perdidas; procesos caídos. &
Cada cliente espera con timeout y, si no hay respuesta, cierra el socket, abre
otro y reenvía la misma petición~\cite{zguide}. Funciona porque las operaciones
son idempotentes por \codigo{idSolicitud}. \\
\textbf{Latidos \emph{Paranoid Pirate}} &
Detectar procesadores caídos sin esperar a que falle una petición. &
Gestor y procesadores intercambian latidos cada 1\,s; tres perdidos sacan al
procesador del grupo. Si el gestor se reinicia, los procesadores se registran de
nuevo solos. \\
\textbf{Broker Majordomo} &
Más servicios y réplicas sin configuración manual. &
El protocolo Majordomo~\cite{rfc7} generaliza el gestor: los procesos se
registran por nombre de servicio (por ejemplo \codigo{rmc}, \codigo{persistencia})
y el broker enruta y balancea. Permite poner también varias réplicas de
persistencia detrás de un broker. \\
\textbf{Par \emph{Binary Star} para el gestor} &
El gestor deja de ser punto único de fallo. &
Un gestor primario y uno de respaldo en máquinas distintas se vigilan con
latidos; si el primario cae, el respaldo toma el control y los clientes
(\emph{Lazy Pirate}) cambian de dirección. El gestor no guarda datos de negocio,
así que el cambio no pierde información. \\
\textbf{Guardar primero y luego publicar + barrido} &
Avisos PUB/SUB perdidos; caída del servicio de cancelación. &
ZeroMQ PUB/SUB no es durable. La solicitud se guarda como \codigo{PENDIENTE}
antes de responder y un barrido cada 2\,s aplica las que no recibieron aviso.
Varias réplicas del servicio de cancelación pueden trabajar a la vez porque
persistencia aplica cada solicitud una sola vez. \\
\textbf{Base de datos replicada (PostgreSQL con CloudNativePG)} &
Persistencia deja de ser punto único de fallo; más capacidad de lectura. &
Una instancia primaria y réplicas con replicación síncrona: si la primaria cae se
promueve una réplica sin perder confirmaciones (RPO = 0). Las consultas pueden
leer de las réplicas y las reservas siempre se deciden en la primaria. La interfaz
\codigo{Almacenamiento} permite cambiar la bitácora por PostgreSQL. \\
\textbf{Kubernetes (1 plano de control + 3 workers)} &
Reinicio automático, reparto de réplicas y escalado. &
Reinicia los procesos caídos (sondas de \emph{liveness}), reparte las réplicas
en nodos distintos (anti-afinidad), reubica los procesos de un nodo caído y
puede escalar los procesadores según la carga (\emph{HorizontalPodAutoscaler}). \\
\textbf{Idempotencia y operaciones atómicas} &
Reintentos sin duplicados; nada aplicado a medias. &
Ya está en el diseño base: cada cambio se valida completo antes de aplicarse y
se guarda con su \codigo{idSolicitud}. Es lo que hace seguros todos los
reintentos anteriores. \\
\textbf{Monitoreo (Prometheus y Grafana)} &
Detectar degradación y fallos. &
Alertas cuando el p95 supera 1\,s, crece la tasa de errores o hay cancelaciones
pendientes antiguas; sirve también para las pruebas de desempeño. \\
\textbf{Seguridad (CurveZMQ y autenticación)} &
Confidencialidad e integridad de los mensajes. &
Cifrado autenticado de todos los sockets y verificación de identidad del
estudiante (sección~\ref{sec:modelos}). \\
\bottomrule
\end{longtable}

\paragraph{Arquitectura escalable para las pruebas.} Para comparar con la base en
las pruebas de la Tabla~1 se usarán las mismas máquinas y programas, cambiando solo
la cantidad de réplicas:

\begin{table}[H]
\centering
\caption{Arquitectura base frente a escalable.}
\begin{tabular}{lcc}
\toprule
Componente & Base & Escalable \\
\midrule
Gestor de Transacciones & 1 & 1 (+ respaldo \emph{Binary Star}) \\
Procesador RMC & 1 & 3 (uno por máquina de servicios) \\
Servicio de cancelación & 1 & 2 \\
Persistencia & 1 & 1 primaria + réplica \\
\bottomrule
\end{tabular}
\end{table}

Se espera que la escalable aumente las transacciones por minuto y reduzca el
tiempo de respuesta cuando hay más clientes, porque en la base todas las
peticiones esperan en la cola de un único procesador RMC.
